\section{Theorem/Lemma Dependency Tracking}

\begin{longtable}{>{\raggedright\arraybackslash}p{2.2cm} >{\raggedright\arraybackslash}p{4.0cm} >{\raggedright\arraybackslash}p{4.2cm} >{\raggedright\arraybackslash}p{4.2cm}}
\toprule
Item & One-line meaning & Depends on & Used later \\
\midrule
\endfirsthead
\toprule
Item & One-line meaning & Depends on & Used later \\
\midrule
\endhead
Lemma 1 & Criterion for spherical-uniform global minimizer in isotropic model & Reduced isotropic risk; model symmetry & Theorem 1 and the example analysis \\
Theorem 1 & Near-optimal SGD in isotropic Gaussian model & Reduced PDE analysis, good initialization, Theorem 3 (explicitly or implicitly) & Main application result \\
Theorem 2 & Near-optimal SGD in anisotropic Gaussian model & Reduced PDE analysis, symmetry reduction, Theorem 3 & Main application result \\
Proposition 1 & $\inf_\theta R_N(\theta)$ is close to $\inf_\rho R(\rho)$ & Mean-field risk representation; moment bound & Theorems 4--5; example interpretation \\
Theorem 3 & Empirical SGD measure converges to DD / diffusion DD & A1--A3; nonlinear dynamics; Supplementary Section 3 & Theorems 1, 2, 5 \\
Proposition 2 & Characterizes fixed points of noiseless DD & DD as a continuity equation; zero-current condition & Theorems 6--7 \\
Proposition 3 & Unique Boltzmann fixed point for noisy DD & Free-energy structure; regularity; diffusion PDE & Theorem 4 \\
Theorem 4 & Diffusion DD converges to unique free-energy minimizer & Proposition 3; A1--A4; dissipation identity \eqref{eq:free-energy-note} (informal reference) & Theorem 5 \\
Theorem 5 & Regularized noisy SGD achieves near-optimal regularized risk & Theorems 3--4; Proposition 1 & Generic algorithmic conclusion \\
Theorem 6 & Positive Hessian implies local stability of atomic fixed point & Proposition 2; local Hessian $H_0$ & Noiseless asymptotic interpretation \\
Theorem 7 & Certain mixed fixed points are unstable for bounded-density initials & Theorem 6-type Hessian geometry; additional level-set assumptions & Noiseless asymptotic interpretation \\
Corollary 2.1 (SI) & Analytic minimizers are either flat or supported on a null set & Analyticity of $V,U$ and minimizer condition & Conceptual explanation for singular minimizers \\
Lemmas 3.1--3.4 (SI) & PDE continuity and trajectory-comparison estimates & A1--A3, concentration bounds & Proof of Theorem 3 \\
\bottomrule
\end{longtable}

\paragraph{Internal vs.\ imported dependencies.}
\begin{itemize}[leftmargin=1.5em]
\item \textbf{Paper-internal backbone:} Proposition~1, Theorem~3, Propositions~2--3, Theorems~4--7.
\item \textbf{Imported background tools:} Wasserstein gradient-flow framework \cite{jordan1998variational,ambrosio2008gradient}; propagation of chaos \cite{sznitman1991topics}.
\item \textbf{Example-specific analysis:} Lemma~1 and the reduced PDEs for isotropic / anisotropic Gaussian models.
\item \textbf{Not part of the proof backbone:} the numerical PDE solutions in the figures. They illustrate the theory, but they do not prove the main abstract convergence statements.
\end{itemize}

\paragraph{Uncertainty marker.}
The paper does not always spell out every proof dependency in theorem-by-theorem style. In particular, the exact dependency of Theorems~1--2 on Theorem~4 is partly inferred from the shared mean-field philosophy and not asserted as a direct proof chain.
